Ein Mol eines Stoffes enthält $N_\text{A}$ Teilchen und hat die Masse $M$ (molare Masse, FoTaBe S.~228--229). Ein Teilchen hat also im Mittel die Masse $m = M/N_\text{A}$.
\begin{abcliste}
\abc Wolfram, $M = \MWoO$:
\begin{align}
 m &= \mWoF = \frac{\MWo}{\ncNA} = \mWo \approx \result{\mWoS{3}}
\end{align}

\abc Sauerstoff, $M = \MOzF = 2\cdot\MOaO = \MOzg$:
\begin{align}
 m &= \mOzF = \frac{\MOz}{\ncNA} = \mOz \approx \result{\mOzS{3}}
\end{align}

\abc Traubenzucker, $M = \MGlF = 6\cdot\MCaO+12\cdot\MHaO+6\cdot\MOaO = \MGlg$:
\begin{align}
 m &= \mGlF = \frac{\MGl}{\ncNA} = \mGl \approx \result{\mGlS{3}}
\end{align}
\end{abcliste}
Die Masse ist eine mittlere Masse, weil die Elemente aus verschiedenen Isotopen bestehen.
